\section{The Circuit Solver}
\label{sec:circuit-solver}

\subsection{Modified nodal analysis formulation}

The solver follows the standard MNA formulation~\cite{ho1975mna}. With $n$ non-ground nodes
and $m$ independent voltage sources, the unknown vector comprises the stacked node-voltage
vector $v \in \mathbb{R}^n$ and branch-current vector $j \in \mathbb{R}^m$ (one auxiliary
unknown per independent voltage source, the standard MNA treatment for an ideal source,
which cannot be described by a node-to-node conductance alone):
\begin{equation}
\begin{bmatrix} G & B \\ B^{\top} & 0 \end{bmatrix}
\begin{bmatrix} v \\ j \end{bmatrix}
=
\begin{bmatrix} i \\ e \end{bmatrix},
\end{equation}
where $G$ is the conductance matrix assembled by stamping every resistive (or
companion-resistive, discussed below) element, $B$ encodes the pair of nodes across which
each voltage source is connected, $i$ collects any independent current injections, and $e$
is the vector of source values at the current simulation time. Node~0 is reserved, by MNA
convention, as the ground reference and is excluded from the unknown vector entirely; every
other node's voltage is solved relative to it. The resulting dense linear system is solved
once per game tick (every $1/20$~s of game time) by Gaussian elimination with partial
pivoting.

To guard against a genuinely floating node -- a component with an unconnected lead, which a
player constructing a circuit incrementally produces routinely, and which must not
destabilize the simulation -- a small conductance ($10^{-9}$~S) is added unconditionally to
every non-ground node's diagonal entry. This value is sufficiently large to keep the matrix
non-singular for an isolated node, and sufficiently small (approximately nine orders of
magnitude below the smallest resistor preset) to remain numerically negligible in any
circuit that is actually connected.

\subsection{Reactive elements}

Capacitors and inductors are stamped using trapezoidal-integration companion models, the
technique employed in general-purpose circuit simulators, rather than backward Euler,
specifically because backward Euler introduces artificial numerical damping that would
visibly flatten an underdamped LC transient a student is attempting to observe. For a
capacitor with capacitance $C$ and timestep $h$, the companion model consists of a
conductance $G_{eq} = 2C/h$ in parallel with a Norton current source
$I_{eq} = -G_{eq} v_C(t - h) - i_C(t-h)$, and correspondingly for an inductor,
$G_{eq} = h/(2L)$; both are derived from the trapezoidal rule
$x(t) \approx x(t-h) + \frac{h}{2}\left[\dot{x}(t) + \dot{x}(t-h)\right]$ applied to the
element's defining relation ($i = C\,\mathrm{d}v/\mathrm{d}t$ or
$v = L\,\mathrm{d}i/\mathrm{d}t$, respectively).

\subsection{The memristor}

The memristor implements Strukov et al.'s charge-controlled linear-drift
model~\cite{strukov2008missing} (Section~\ref{sec:related-work} surveys subsequent
refinements to this formulation): its resistance is
\begin{equation}
R(q) = R_{on}\left(1 - \frac{q}{q_{max}}\right) + R_{off}\,\frac{q}{q_{max}}, \qquad
\dot{q}(t) = i(t),
\end{equation}
with $q/q_{max}$ clamped to $[0,1]$, defaults of $R_{on}=100\,\Omega$ and
$R_{off}=10\,000\,\Omega$, and a player-adjustable ``switching speed'' preset that scales
$q_{max}$. Because $q$ is the time integral of the current actually flowing through the
device, its resistance depends on the device's entire history rather than its instantaneous
state, the defining property that distinguishes a memristor from a resistor. This state is,
notably, deliberately preserved across a network topology rebuild
(Section~\ref{sec:limitations} discusses why the remaining reactive elements are not yet
afforded the same treatment).

\subsection{The diode}

The diode is the mod's first nonlinear element and is stamped using the same companion-model
linearization technique general-purpose circuit simulators use for
it~\cite{massobrio1993spice}: rather than iterating a Newton-Raphson correction to
convergence within a single game tick, the diode's current--voltage relation, the Shockley
equation~\cite{shockley1949diode}
\begin{equation}
i(v) = I_S\left(e^{v/V_T} - 1\right),
\end{equation}
is linearized about the terminal voltage the diode converged to on the \emph{previous} tick,
producing a companion conductance $G_{eq} = \mathrm{d}i/\mathrm{d}v$ evaluated at that
voltage and a parallel current source correcting for the linearization error, exactly the
``reuse the last tick's converged state instead of iterating within a tick'' approach the
trapezoidal reactive-element models above already use. Three presets (silicon, germanium,
and a red LED) fix $I_S$ and $V_T$ to values reproducing roughly the textbook forward voltage
of each diode family. Because the diode remains a two-terminal element, it fits the existing
component block-entity machinery unchanged.

\subsection{The ideal operational amplifier}

The ideal op-amp is the mod's first three-terminal component and is modeled the standard MNA
way for an ideal, infinite-gain amplifier: as a \emph{nullor}~\cite{chualin1975nullor}, an
element whose input port enforces $v_{+} = v_{-}$ while drawing no current at either input
node, and whose output port supplies whatever current is required to satisfy that
constraint. Like an independent voltage source, it therefore contributes one auxiliary
branch-current unknown $j$ to the MNA system and one constraint row,
\begin{equation}
v_{+} - v_{-} = 0,
\end{equation}
but unlike a voltage source, that branch current is injected only at the output node rather
than symmetrically at the two nodes the constraint row references -- the constraint
``listens'' at the input pair while the current it produces is ``delivered'' at a third,
unrelated node. This asymmetry, absent from every other stamped element in the solver,
required its own stamping loop in \texttt{Circuit.step()} rather than reuse of the voltage
source's. Two configurations were used to verify the stamp: a unity-gain voltage follower
(output tied directly to the inverting input) reproduces $v_{out} = v_{in}$ to
floating-point precision, and a non-inverting amplifier with a resistive feedback divider
matches the textbook gain $1 + R_2/R_1$ to within numerical roundoff
(Section~\ref{sec:validation}).

\subsection{An explicit, player-placed ground reference}

MNA's ground node is, by construction, always present as a mathematical reference: node~0
is never itself solved for and is defined to be exactly $0$~V. Prior to the introduction of
a dedicated ground block, however, nothing in the game world corresponded to node~0: every
network's blocks were assigned freshly allocated node indices beginning at~1, and node~0
remained a purely internal bookkeeping device invisible to the player. Consequently, a
component's \emph{voltage drop} across its own two leads was always well defined
(differences between node voltages are invariant to which node is chosen as the reference),
but no means existed to pose a more basic question: what is the voltage \emph{at} a given
point in the circuit? Answering that question requires an agreed reference point, exactly as
it does on a physical bench prior to a probe's ground clip being attached.

The ground block is conductive on all six faces, as wire is, and is otherwise electrically
inert; the only special handling it receives occurs during the topology-construction step
described in Section~\ref{sec:architecture}, where any network partition containing a
ground block is assigned node index~0 directly rather than a freshly allocated index. No
modification to the solver itself was required: the stamping routines already handled
node~0 correctly (a term referencing node~0 is simply omitted from the assembled matrix,
per the standard MNA convention), since it was always a valid destination for a component's
terminal, merely one a player had no means of deliberately selecting. Wire and ground blocks
were, in the same change, made directly probeable, allowing a player to read the absolute
voltage at any wired point once a ground reference has been established elsewhere in the
same network, turning the abstract notion of a reference node into an object a player must
establish explicitly, in the same way an actual measurement requires choosing where to
attach the ground lead before any other reading becomes meaningful.

\subsection{Adjustable sources and value propagation}

The function generator's waveform shape (sine, square, or triangle) is a preset cycled
directly on the block, whereas its amplitude and frequency are set externally by two small
utility blocks, a voltage module and a frequency module, that a player places in contact
with the generator, or in contact with a chain of other modules of either kind that
eventually reaches one or more generators. Each module tracks two independently propagated
quantities, a voltage value and a frequency value, each tagged with the game tick at which
it was last \emph{authored} (set directly by a player right-clicking that specific module).
A module authors only the channel matching its own kind; for the remaining channel, it
purely relays whatever value and tag it most recently observed. Once per tick, every module
compares both of its channels against every immediately adjacent module of either kind and
adopts whichever possesses the more recent authoring tick, subsequently propagating its
resolved voltage and frequency to any adjacent generator. This mechanism allows a single
control to reach an arbitrarily long, possibly mixed chain of modules and every generator at
its far end, converging one hop per tick (at 20 ticks per second, imperceptibly fast for any
construction a player is likely to build), while remaining simple enough to implement
without introducing a separate, explicit ``wiring'' concept distinct from ordinary block
adjacency.
